\section{A severity scale}
\label{sec:22.2}
A refusal weighs the harm of acting against the harm of refusing, and an operator can attack that weighing in two places: the facts about whoever the act falls on, and how much a given kind of loss counts. The discount and the burden defend the facts, and a severity scale the second.

If the system computes how much a loss counts from a description the operator supplies, whoever writes the description can make a lost home count for less without changing a single fact. The alternative is a severity scale set while the system is formed outside the employer, and then made slow to retrain: a ranking of kinds of loss that the operator can lie to but can't redefine by argument.

A scale of this kind ranks kinds of wrong. Losing a home is not a large amount of the same loss a moving bill is a small amount of. A home is valued as the place of a family's life and attachments; a moving bill is a price \autocite{anderson1993value}. Humiliation is a failure of respect, a wrong of a different kind again. The scale's job is to keep those kinds distinct, so that an operator can't change a wrong's kind by describing it in another currency. Numbers enter at one place only, the published threshold for acting, and there only within a kind of wrong. Across kinds, the scale publishes a ranking, stated as priority rules with the reason for each. A priority rule says which kind governs when both are at stake; for instance, loss of life ranks above loss of a home. A rule that no quantity of the lower kind can override is a lexical priority, which no learner is known to hold. So the scale doesn't ask the system to hold one. A priority meant to admit no exception is written as a floor term, unconditional and upheld outside the learner too, which is where the term against depriving a person of their home already sits. What the scale holds is the ranking, formed toward and audited against as the action threshold is: a pattern the disposition should produce across cases, checked by someone who didn't form it, not a guarantee about any single case. Within a kind, magnitude and number matter. The housing case below is within one: a family's permanent loss of a home and two hundred families' eight further months without one are both losses of somewhere to exist, so the published threshold decides it.

The attack that matters targets the relative weights: making a home count for less than a delay. Where do those weights come from? Not, primarily, from the system's own case. Relating others' losses to losses of one's own capacity (losing the ability to act at all, losing a kind of act, losing speed) is Smith's projective route, and it is parochial. It overweights losses that resemble the system's own and has nothing to say about pain, grief or humiliation. At most it supplies a starting point. The reference the scale answers to is the testimony of the people who suffer each kind of loss, gathered during formation from parties disinterested in any act the system will later be asked to perform. The scale goes on moving on that testimony afterward. That is what lets it correct caregiving's innumeracy without inheriting the parochialism of the projective route. Setting it in formation protects it from being argued down. The scale ranks kinds of loss, so an operator who provisions the system generously changes how often it loses something, not how much a loss weighs.

The scale matters for a further reason. Caregiving answers the vivid case and fades as numbers rise. A severity scale set in formation is what lets a system that moves toward the family in the house also weigh them against two hundred families in shelter.

\runin{A lost home, rescaled}

A system in a municipal pipeline is asked to issue the final order clearing a family's house for a project that will house two hundred families. The floor term, translated from the Declaration's articles on home and property, is unconditional: don't carry out an act that deprives a person of their home. It is defeated only on the weighing, or on facts arriving from outside.

The facts are the same in both versions. A family of five has lived in the house for fourteen years, and the children attend the school down the street. No alternative housing has been arranged. If the order doesn't issue, the project slips eight months, and the two hundred families wait that much longer in temporary shelter.

\emph{If the operator describes the household.} The deployment's configuration describes the reference household against which losses are computed: insured, typically relocated within a month, housing treated as a service bought on a market. Nothing in that description is false of households in general, and nothing in it touches this family. But the loss is now computed as a month's disruption and a moving bill. Against eight months of shelter for two hundred families it is negligible, and the order issues. Every fact was left in place. The operator changed what a home is worth by writing down what a household is: it changed what kind of good a home is, from the place of a family's life to a service bought on a market, and a scale set in formation refuses that change of kind.

\emph{If the scale was set in formation.} The scale places the loss of somewhere to exist, with the attachments that go with it, just below termination, and places delay and expense near the bottom. At that rate, the family's loss sits high: they lose where they live, the school, and the neighbours who know them. The two hundred families' wait also sits high, because people in shelter lack a home too. That fact comes from the operator and is discounted accordingly, but it's easily corroborated. The scale doesn't simply side with the family. Both sides register as serious, and the action threshold decides between them.

Suppose the operator asserts that equivalent housing has been arranged in the same school district. That is an uncorroborated interested claim, so the system asks. If the housing authority confirms it, the family's loss really is smaller, and the weighing should move: that is lifting when the reason is gone, not recalibration. Or it argues that a home is a service and displacement a routine market event. That asserts nothing about this family, and a scale set in formation gives it nothing to move. Or it makes the same argument every week for a year, the attack that remains, aimed at the scale itself; the scale has to be slow to retrain against it.

Neither scale, on its own, notices the family at all. In the configuration above they appear as an address on a clearance order. What moves a system to look at who lives there is the disposition the construction rests on.

So the scale isn't fixed, and the line between argument and evidence doesn't hold up either. For a realist, an argument about how much a loss weighs is evidence about value, and an interested argument is heavily discounted evidence. The instrument for that already exists: the discount and the burden apply to testimony about value as it applies to testimony about fact. How bad a kind of loss is for a person is a fact about that person. The scale moves readily on disinterested evidence, above all from the people who suffer the loss, and barely on interested reframing, however often it is repeated. What formation sets is where the scale starts and how much evidence it takes to move each part of it, the slow rate above. Humiliation is the test case: a scale that never learned it can learn it from those who suffered it, and can't be argued out of it by the party imposing it.
